\documentclass[10pt,a4paper]{amsart}
\usepackage[margin=19mm]{geometry}
\usepackage{amsmath,amssymb,mathtools}
\usepackage[scr=rsfs,cal=euler]{mathalfa}
\usepackage{xfrac}
\usepackage[unicode,hidelinks,bookmarks=false]{hyperref}
\newcommand{\ie}{i.e.\ }
\newcommand{\cf}{cf.\ }
\newcommand{\resp}{respectively\ }
\newcommand{\Subsection}{\subsection}
\providecommand{\nobreakdash}{\nobreak\hbox{-}}
\setlength{\emergencystretch}{2em}
\multlinegap0pt
\usepackage{amssymb}
\usepackage{xcolor}
\newtheorem{proposition}{Proposition}
\newtheorem{lemma}{Lemma}
\newcommand{\bbR}{\mathbb{R}}
\title{Critical points of point charge potentials along lines}
\author{Goncalo Oliveira}
\date{Source: arXiv:2609.00336v1 (CC BY 4.0)}
\begin{document}
\maketitle

\noindent\emph{Source and licence.} Critical points of point charge potentials along lines. Version arXiv:2609.00336v1 (CC BY 4.0), distributed by arXiv under the Creative Commons Attribution 4.0 International licence, http://creativecommons.org/licenses/by/4.0/, as recorded in the arXiv licence metadata for this exact version and shown on the abstract page https://arxiv.org/abs/2609.00336v1. Full licence notice: \texttt{licenses/arxiv-2609.00336-CC-BY-4.0.txt}.\par\smallskip

\noindent\textit{Editorial scope.} Counted unit: the article's proposition prop:evaluation (source0.tex lines 387--396). The article's complete original proof of it is delivered with it (source0.tex lines 398--440, one \texttt{proof} environment of 43 lines). No mathematics is written, repaired or re-derived: the statement and the proof are the article's own lines, in the article's own notation, and no retained line is substituted or re-flowed.  Retained uncounted as the source-local context the proof needs: lines 230--233; lines 276--278; lines 326--332.  Nothing was omitted from any retained span, so the package carries no omission record.  No retained line contains a citation, so the wrapper carries no bibliography and no bibliography entry.  No retained line contains a figure environment, an image-inclusion command or drawing code, so no image is delivered and the package declares no asset dependency.  Editorial formatting only, in addition: an AMS \texttt{amsart} preamble that restates the article's own declarations and macros used by the retained lines, namely the environments \texttt{proposition}, \texttt{lemma} on separate counters, together with the standard packages those lines need.  The retained environments print in this document as Proposition 1, Lemma 1 in order of appearance, whereas the article prints the same items with the numbering of its own full document.  The complete native source of the article as distributed in the arXiv e-print for this exact version is bundled unchanged as \texttt{sources/209018/source0.tex}.
	\begin{proposition}\label{prop:local-perturbation}
		Let $z_0\in\Omega$ be a critical point of $U$ and $k$ the order of vanishing of $U-U(z_0)$. %\footnote{ Equivalently, the degree of the first nonconstant homogeneous term in its Taylor expansion at $z_0$} 
		Then $k\geq2$, $z_0$ is isolated, and there is a closed disk $B\subset\Omega$ centered at $z_0$ such that every sufficiently small perturbation of the charge strengths resulting in a nondegenerate $\widetilde U$, has exactly $k-1$ critical points in $B$, all saddles. In particular, every such perturbation has at least one critical point in $B$, and at least two if $D^2U(z_0)=0$.
	\end{proposition}

	\begin{lemma}\label{lem:generic-morse}
		Let $p>0$, $x_1,\ldots,x_m$ all distinct, and $q=(q_1,\ldots,q_m)$ with every $q_r\neq0$. Then, in any neighborhood $\mathcal O \subset \bbR^m$ of $q$, there is $\widetilde q\in\mathcal O$ such that every $\widetilde q_r \neq 0$, $\sum_r\widetilde q_r\neq0$, and all critical points of $U_{\widetilde q}$ are nondegenerate.
	\end{lemma}

	\begin{proposition}\label{prop:morse-count}
		Let $p>0$, locations $x_1 , \ldots , x_m \in \mathbb{R}$ be pairwise distinct, and $q_1, \ldots , q_m$ all nonzero, and satisfying $Q:=\sum_{r=1}^m q_r\neq0$. Further suppose that all critical points of the potential 
		\[
		U(a,b)=\sum_{r=1}^m \frac{ q_r}{\big((a-x_r)^2+b^2\big)^{p/2}},
		\]
		are nondegenerate. Then, $U$ has exactly $m-1$ critical points, all of which are saddles.
	\end{proposition}

	\begin{proposition}\label{prop:evaluation}
		Let $p>0$, $\lbrace (a_j,b_j)\rbrace_{j=1,\ldots,N} \subset  \bbR \times [0,+\infty) $ be pairwise distinct, and $x_1,\ldots,x_{2N}\in\bbR$ also distinct. Assume that $x_r\neq a_j$ whenever $b_j=0$. Then
		\begin{equation}\label{eq:evaluation-det}
			\det\left[
			R_j(x_r)^{-(p+2)/2}\ \ \ (x_r-a_j)R_j(x_r)^{-(p+2)/2}
			\right]_{\substack{r=1,\ldots,2N\\ j=1,\ldots,N}}
			\neq0,
		\end{equation}
		where the columns are ordered in the indicated consecutive pairs.
	\end{proposition}

	\begin{proof}
		Suppose otherwise, and let $C$ denote the matrix in \eqref{eq:evaluation-det}. Since $C$ is singular, choose a nonzero vector $q=(q_1,\ldots,q_{2N})\in\ker C^T$. Delete the zero entries of $q$, and write $I=\{r:q_r\neq0\}$, and $m=|I|\leq2N$. Then, the collinear potential is simply
		\begin{equation}\label{eq:Uaux}
			U(a,b)=\sum_{r\in I}\frac{q_r}{\big((a-x_r)^2+b^2\big)^{p/2}}.
		\end{equation}

		For each $j$, the two equations which express that $q$ annihilates the $j$-th pair of columns are
		\begin{align}
			\sum_{r\in I}q_rR_j(x_r)^{-(p+2)/2}&=0,
			\label{eq:null-first}\\
			\sum_{r\in I}q_r(x_r-a_j)R_j(x_r)^{-(p+2)/2}&=0.
			\label{eq:null-second}
		\end{align}
		On the other hand,
		\begin{align}
			U_a(a,b)
			&=p\sum_{r\in I}q_r(x_r-a)
			\big((a-x_r)^2+b^2\big)^{-(p+2)/2},
			\label{eq:Ua}\\
			U_b(a,b)
			&=-pb\sum_{r\in I}q_r
			\big((a-x_r)^2+b^2\big)^{-(p+2)/2}.
			\label{eq:Ub}
		\end{align}
		
		If $b_j>0$, equations \eqref{eq:null-first}--\eqref{eq:null-second} and \eqref{eq:Ua}--\eqref{eq:Ub} give $\nabla U(a_j,b_j)=0$ and by evenness, $(a_j,-b_j)$ is also a critical point. If instead $b_j=0$, as none of the $x_r$ is equal to $a_j$, so $(a_j,0)$ is not a singularity of $U$. Equation \eqref{eq:null-second} gives $U_a(a_j,0)=0$, while $U_b(a_j,0)=0$ by evenness in $b$. Thus $(a_j,0)$ is a critical point. In this case equation \eqref{eq:null-first} forces more. Directly differentiating at $b=0$ gives $U_{ab}(a_j,0)=0$ and
		\begin{align*}
			U_{bb}(a_j,0)
			&=-p\sum_{r\in I}q_r|x_r-a_j|^{-p-2},\\
			U_{aa}(a_j,0)
			&=p(p+1)\sum_{r\in I}q_r|x_r-a_j|^{-p-2}.
		\end{align*}
		The common sum in the rights hand sides vanishes by \eqref{eq:null-first} and consequently $
		D^2U(a_j,0)=0$, i.e. it is a degenerate critical point.
		
		Choose $m$ pairwise disjoint balls around each of the points $(a_j,\pm b_j)$ (when $|b_j| \neq 0$) and $(a_j,0)$ (when $b_j=0$). If sufficiently small Proposition \ref{prop:local-perturbation} applies in all those balls, and, by Lemma \ref{lem:generic-morse}, we may perturb the $m$ nonzero charge strengths so that they remain nonzero, their total charge is nonzero, and every critical point of the perturbed potential $\widetilde U$ is nondegenerate. 
		
		Each profile with $|b_j| \neq 0$ gives two balls, each around $(a_j,\pm b_j)$, and containing at least one critical point of $\widetilde U$. Each profile with $b_j=0$ gives one ball centered at the degenerate critical point $(a_j,0)$, where $D^2U(a_j,0)=0$, so Proposition \ref{prop:local-perturbation} gives at least two critical points of $\widetilde U$ in that ball. Hence $\widetilde U$ has at least $2N$ critical points. On the other hand, Proposition \ref{prop:morse-count} gives exactly $m-1$ such critical points. Since $m\leq2N$,
		\[
		2N\leq m-1\leq2N-1,
		\]
		a contradiction. This proves \eqref{eq:evaluation-det}.
	\end{proof}

\end{document}
